\documentclass[11pt]{article}

\usepackage[a4paper,margin=1in]{geometry}
\usepackage{amsmath,amssymb,amsthm,mathtools}
\usepackage{enumitem}
\usepackage{microtype}
\usepackage[hidelinks]{hyperref}
\usepackage{algorithm}
\usepackage{algpseudocode}

\newtheorem{remark}{Remark}
\newcommand{\R}{\mathbb{R}}
\newcommand{\N}{\mathbb{N}}
\newcommand{\abs}[1]{\left\lvert #1\right\rvert}
\newcommand{\norm}[1]{\left\lVert #1\right\rVert}

\title{Certified Integration over Polynomial-Zonotope Domains}
\author{}
\date{}

\begin{document}
\maketitle

\section{Setting and interpretation}
\label{sec:pz-integration}

Let
\[
\Xi=[-1,1]^s
\]
be a reference parameter box, and let
\[
X:\Xi\to\R^n
\]
be a polynomial map of the form
\begin{equation}
    X(\alpha)
    =
    c+\sum_{\beta\in\mathcal{B}} g_\beta \alpha^\beta,
    \qquad
    \alpha=(\alpha_1,\ldots,\alpha_s)\in\Xi,
    \label{eq:pz-parametrization}
\end{equation}
where
\[
    c,g_\beta\in\R^n,
    \qquad
    \alpha^\beta
    :=
    \prod_{j=1}^s \alpha_j^{\beta_j}.
\]
The corresponding polynomial zonotope is
\[
    \mathcal{X}=X(\Xi).
\]

Let \(f:\mathcal{X}\to\R\) be a scalar-valued function. Assume that a certified polynomial enclosure of \(f\) on \(\mathcal{X}\) has been computed:
\begin{equation}
    f(x)=p(x)+r(x),
    \qquad
    \abs{r(x)}\leq E
    \quad\text{for all }x\in\mathcal{X},
    \label{eq:polynomial-remainder-enclosure}
\end{equation}
where \(p:\R^n\to\R\) is a polynomial and \(E\geq 0\).

Equivalently, one may write the pointwise enclosure as
\begin{equation}
    f(x)\in p(x)+E[-1,1].
    \label{eq:pointwise-noise}
\end{equation}
The interval in \eqref{eq:pointwise-noise} is a pointwise remainder enclosure. In general, it does not mean that there exists one constant noise value \(\eta\in[-1,1]\) such that
\[
    f(x)=p(x)+E\eta
\]
simultaneously for all \(x\in\mathcal{X}\). Rather, the remainder is an unknown function \(r(x)\) satisfying \(\abs{r(x)}\leq E\).

There are two different integrals associated with the parametrization \eqref{eq:pz-parametrization}:
\begin{enumerate}[label=\arabic*.]
    \item the parameter-space integral
    \[
        \int_\Xi f(X(\alpha))\,d\alpha,
    \]
    \item the geometric integral
    \[
        \int_{\mathcal{X}} f(x)\,dx,
    \]
    where \(dx\) denotes \(n\)-dimensional Lebesgue measure.
\end{enumerate}
These integrals agree only in special cases. The geometric integral requires an appropriate Jacobian factor.

\section{Exact integration of polynomials on the reference box}

Let
\[
    q(\alpha)=\sum_{\gamma\in\mathcal{G}} q_\gamma\alpha^\gamma
\]
be a polynomial on \(\Xi=[-1,1]^s\). Then
\begin{equation}
    \int_\Xi q(\alpha)\,d\alpha
    =
    \sum_{\gamma\in\mathcal{G}} q_\gamma\,\mu_\gamma,
    \label{eq:box-polynomial-integration}
\end{equation}
where
\begin{equation}
    \mu_\gamma
    :=
    \int_\Xi \alpha^\gamma\,d\alpha
    =
    \prod_{j=1}^s\int_{-1}^1 t^{\gamma_j}\,dt.
    \label{eq:box-moment}
\end{equation}
For every \(k\in\N_0\),
\begin{equation}
    \int_{-1}^1 t^k\,dt
    =
    \begin{cases}
        0, & k\text{ odd},\\[1mm]
        \displaystyle\frac{2}{k+1}, & k\text{ even}.
    \end{cases}
    \label{eq:one-dimensional-moment}
\end{equation}
Consequently,
\begin{equation}
    \mu_\gamma
    =
    \begin{cases}
        \displaystyle\prod_{j=1}^s \frac{2}{\gamma_j+1},
        & \gamma_j\text{ is even for every }j,\\[3mm]
        0, & \text{otherwise}.
    \end{cases}
    \label{eq:multivariate-moment}
\end{equation}
Thus polynomial integration over the reference box is a linear coefficient operation. No range enclosure of the polynomial is needed.

\section{Integration with respect to parameter measure}

Define the polynomial pullback
\begin{equation}
    q(\alpha):=(p\circ X)(\alpha).
    \label{eq:pullback-polynomial}
\end{equation}
Since both \(p\) and \(X\) are polynomial, \(q\) is polynomial.

From \eqref{eq:polynomial-remainder-enclosure},
\[
    f(X(\alpha))=q(\alpha)+r(X(\alpha)),
    \qquad
    \abs{r(X(\alpha))}\leq E.
\]
Therefore,
\begin{align}
    \int_\Xi f(X(\alpha))\,d\alpha
    &=
    \int_\Xi q(\alpha)\,d\alpha
    +
    \int_\Xi r(X(\alpha))\,d\alpha \\
    &\in
    \int_\Xi q(\alpha)\,d\alpha
    +
    E\,\abs{\Xi}[-1,1].
\end{align}
Since \(\abs{\Xi}=2^s\),
\begin{equation}
    \boxed{
    \int_\Xi f(X(\alpha))\,d\alpha
    \in
    I_{\mathrm{poly}}+2^sE[-1,1],
    }
    \label{eq:parameter-integral-enclosure}
\end{equation}
where
\[
    I_{\mathrm{poly}}:=\int_\Xi q(\alpha)\,d\alpha.
\]
Equivalently, introducing one fresh noise symbol,
\begin{equation}
    \int_\Xi f(X(\alpha))\,d\alpha
    \in
    I_{\mathrm{poly}}+2^sE\,\eta_{\mathrm{int}},
    \qquad
    \eta_{\mathrm{int}}\in[-1,1].
    \label{eq:parameter-integral-noise}
\end{equation}

\section{Geometric integration over the polynomial-zonotope image}

Assume now that \(s=n\), and that \(X:\Xi\to\mathcal{X}\subset\R^n\) is injective and continuously differentiable, with
\[
    \det DX(\alpha)\neq 0
    \qquad
    \text{for all }\alpha\in\Xi.
\]
Then the change-of-variables formula gives
\begin{equation}
    \int_{\mathcal{X}} f(x)\,dx
    =
    \int_\Xi f(X(\alpha))J_X(\alpha)\,d\alpha,
    \label{eq:change-of-variables}
\end{equation}
where
\begin{equation}
    J_X(\alpha):=\abs{\det DX(\alpha)}.
    \label{eq:jacobian-density}
\end{equation}
Using \eqref{eq:polynomial-remainder-enclosure},
\begin{align}
    \int_{\mathcal{X}} f(x)\,dx
    &=
    \int_\Xi p(X(\alpha))J_X(\alpha)\,d\alpha
    +
    \int_\Xi r(X(\alpha))J_X(\alpha)\,d\alpha.
\end{align}
Since \(J_X(\alpha)\geq 0\),
\[
    \abs{\int_\Xi r(X(\alpha))J_X(\alpha)\,d\alpha}
    \leq
    E\int_\Xi J_X(\alpha)\,d\alpha.
\]
Moreover,
\begin{equation}
    \abs{\mathcal{X}}=\int_\Xi J_X(\alpha)\,d\alpha.
    \label{eq:image-volume}
\end{equation}
Hence
\begin{equation}
    \boxed{
    \int_{\mathcal{X}} f(x)\,dx
    \in
    I_{\mathrm{poly}}+E\abs{\mathcal{X}}[-1,1],
    }
    \label{eq:geometric-integral-general}
\end{equation}
where
\begin{equation}
    I_{\mathrm{poly}}
    :=
    \int_\Xi p(X(\alpha))J_X(\alpha)\,d\alpha.
    \label{eq:geometric-polynomial-part}
\end{equation}

\section{Fixed-orientation case}

Suppose that the sign of the Jacobian determinant is known:
\begin{equation}
    \sigma\det DX(\alpha)>0
    \qquad
    \text{for all }\alpha\in\Xi,
    \qquad
    \sigma\in\{-1,1\}.
    \label{eq:fixed-orientation}
\end{equation}
Then
\[
    J_X(\alpha)=\sigma\det DX(\alpha).
\]
Because \(X\) is polynomial, both \(\det DX(\alpha)\) and \(p(X(\alpha))\det DX(\alpha)\) are polynomials in \(\alpha\). Therefore,
\begin{equation}
    I_{\mathrm{poly}}
    =
    \sigma\int_\Xi p(X(\alpha))\det DX(\alpha)\,d\alpha,
    \label{eq:fixed-orientation-poly-integral}
\end{equation}
and
\begin{equation}
    \abs{\mathcal{X}}
    =
    \sigma\int_\Xi \det DX(\alpha)\,d\alpha
    \label{eq:fixed-orientation-volume}
\end{equation}
can both be computed by exact polynomial moment integration.

\section{Affine-zonotope special case}

Suppose
\begin{equation}
    X(\alpha)=c+G\alpha,
    \qquad
    G\in\R^{n\times n},
    \qquad
    \det G\neq 0.
    \label{eq:affine-zonotope-map}
\end{equation}
Then \(DX(\alpha)=G\) and \(J_X(\alpha)=\abs{\det G}\) is constant. Hence
\begin{equation}
    \int_{\mathcal{X}} f(x)\,dx
    \in
    \abs{\det G}\int_\Xi p(c+G\alpha)\,d\alpha
    +
    E\abs{\det G}\,2^n[-1,1].
    \label{eq:affine-zonotope-integral}
\end{equation}
Since \(\abs{\mathcal{X}}=2^n\abs{\det G}\), this becomes
\begin{equation}
    \boxed{
    \int_{\mathcal{X}} f(x)\,dx
    \in
    \abs{\det G}\int_\Xi p(c+G\alpha)\,d\alpha
    +
    E\abs{\mathcal{X}}[-1,1].
    }
    \label{eq:affine-zonotope-boxed}
\end{equation}

\section{Jacobian sign changes and injectivity}

If \(\det DX\) changes sign on \(\Xi\), then \(\abs{\det DX}\) is generally only piecewise polynomial. A practical certified procedure is to subdivide
\[
    \Xi=\bigcup_{\ell=1}^N \Xi_\ell
\]
into boxes with pairwise disjoint interiors until the sign of \(\det DX\) can be certified on every \(\Xi_\ell\). On each subbox, choose \(\sigma_\ell\in\{-1,1\}\) such that
\[
    \sigma_\ell\det DX(\alpha)>0
    \qquad
    \text{for all }\alpha\in\Xi_\ell.
\]
Then each local contribution is again the integral of a polynomial after affine rescaling to \([-1,1]^n\).

If injectivity of \(X\) is not known, the parameter-space integral
\[
    \int_\Xi f(X(\alpha))\abs{\det DX(\alpha)}\,d\alpha
\]
may count points of the image with multiplicity. Therefore certified geometric integration over \(\mathcal{X}\) requires either a proof that \(X\) is injective on \(\Xi\), or a subdivision into pieces on which \(X\) is injective and whose images overlap only on sets of measure zero.

\section{Spatially varying interval remainder}

Assume more generally that
\begin{equation}
    f(x)=p(x)+r(x),
    \qquad
    \abs{r(x)}\leq \rho(x)
    \quad
    \text{for all }x\in\mathcal{X},
    \label{eq:variable-remainder}
\end{equation}
where \(\rho:\mathcal{X}\to[0,\infty)\). Then
\begin{equation}
    \int_{\mathcal{X}} f(x)\,dx
    \in
    \int_{\mathcal{X}}p(x)\,dx
    +
    \left[
        -\int_{\mathcal{X}}\rho(x)\,dx,
        \int_{\mathcal{X}}\rho(x)\,dx
    \right].
    \label{eq:variable-remainder-integral}
\end{equation}
If \(\rho\circ X\) is polynomial and the Jacobian sign is fixed, then the remainder radius can itself be computed exactly by polynomial moment integration.

\section{Several cells}

Suppose
\[
    \Omega=\bigcup_{k=1}^N\mathcal{X}_k
\]
with pairwise disjoint interiors, and assume
\[
    f(x)=p_k(x)+r_k(x),
    \qquad
    \abs{r_k(x)}\leq E_k
    \quad
    \text{for }x\in\mathcal{X}_k.
\]
Let
\[
    I_k^{\mathrm{poly}}:=\int_{\mathcal{X}_k}p_k(x)\,dx,
    \qquad
    V_k:=\abs{\mathcal{X}_k}.
\]
Then
\begin{equation}
    \boxed{
    \int_\Omega f(x)\,dx
    \in
    \sum_{k=1}^N I_k^{\mathrm{poly}}
    +
    \left(\sum_{k=1}^N E_kV_k\right)[-1,1].
    }
    \label{eq:multi-cell-integral}
\end{equation}
For a scalar final integral, a single fresh final noise symbol is sufficient:
\begin{equation}
    \int_\Omega f(x)\,dx
    \in
    \sum_{k=1}^N I_k^{\mathrm{poly}}
    +
    \left(\sum_{k=1}^N E_kV_k\right)\eta_{\mathrm{int}},
    \qquad
    \eta_{\mathrm{int}}\in[-1,1].
\end{equation}

\section{Coefficient-level implementation}

Let
\[
    h(\alpha)=\sum_{\gamma\in\mathcal{G}} h_\gamma\alpha^\gamma.
\]
Define the integration operator
\begin{equation}
    \mathcal{I}(h)
    :=
    \sum_{\gamma\in\mathcal{G}} h_\gamma\mu_\gamma,
    \label{eq:integration-operator}
\end{equation}
where \(\mu_\gamma\) is given by \eqref{eq:multivariate-moment}. The operator \(\mathcal{I}\) is linear, so it should be applied directly to sparse polynomial coefficients without first range-enclosing the polynomial.

\begin{algorithm}[ht]
\caption{Exact polynomial integration on \([-1,1]^s\)}
\label{alg:polynomial-box-integration}
\begin{algorithmic}[1]
\Require Sparse polynomial \(h(\alpha)=\sum_{\gamma\in\mathcal{G}}h_\gamma\alpha^\gamma\)
\Ensure \(I=\int_{[-1,1]^s}h(\alpha)\,d\alpha\)
\State \(I\gets 0\)
\ForAll{\(\gamma\in\mathcal{G}\)}
    \If{some component \(\gamma_j\) is odd}
        \State \(\mu_\gamma\gets 0\)
    \Else
        \State \(\mu_\gamma\gets\prod_{j=1}^s\frac{2}{\gamma_j+1}\)
    \EndIf
    \State \(I\gets I+h_\gamma\mu_\gamma\)
\EndFor
\State \Return \(I\)
\end{algorithmic}
\end{algorithm}

\begin{algorithm}[ht]
\caption{Certified integration over a polynomial-zonotope image}
\label{alg:certified-pz-integration}
\begin{algorithmic}[1]
\Require Polynomial parametrization \(X:\Xi=[-1,1]^n\to\R^n\)
\Require Polynomial \(p:\R^n\to\R\)
\Require Certified bound \(\abs{f(x)-p(x)}\leq E\) on \(\mathcal{X}=X(\Xi)\)
\Require Certificate that \(X\) is injective on \(\Xi\)
\Require Certificate \(\sigma\det DX(\alpha)>0\) on \(\Xi\)
\State \(d(\alpha)\gets\sigma\det DX(\alpha)\)
\State \(q(\alpha)\gets p(X(\alpha))\)
\State \(h(\alpha)\gets q(\alpha)d(\alpha)\)
\State \(I_{\mathrm{poly}}\gets\int_\Xi h(\alpha)\,d\alpha\)
\State \(V\gets\int_\Xi d(\alpha)\,d\alpha\)
\State \(R\gets EV\)
\State \Return \([I_{\mathrm{poly}}-R,\,I_{\mathrm{poly}}+R]\)
\end{algorithmic}
\end{algorithm}

\section{Floating-point rigor}

The mathematical formulas above are exact. A computer implementation must additionally enclose floating-point roundoff. Two common approaches are:
\begin{enumerate}[label=\arabic*.]
    \item store the polynomial coefficients as intervals and perform all coefficient operations with outward rounding;
    \item store floating-point coefficients and compute an additional certified roundoff remainder.
\end{enumerate}
If
\[
    I_{\mathrm{poly}}
    \in
    [\underline{I}_{\mathrm{poly}},\overline{I}_{\mathrm{poly}}]
\]
and
\[
    \abs{\mathcal{X}}
    \in
    [\underline{V},\overline{V}],
    \qquad
    0\leq\underline{V}\leq\overline{V},
\]
then a rigorous final enclosure is
\begin{equation}
    \int_{\mathcal{X}}f(x)\,dx
    \in
    [\underline{I}_{\mathrm{poly}}-E\overline{V},
     \overline{I}_{\mathrm{poly}}+E\overline{V}].
    \label{eq:roundoff-safe-final}
\end{equation}

\section{Recommended computational pipeline}

The recommended procedure is
\[
\boxed{
\begin{aligned}
&\text{polynomial-zonotope parametrization}
\\
&\quad\longrightarrow\text{polynomial pullback}
\\
&\quad\longrightarrow\text{Jacobian multiplication}
\\
&\quad\longrightarrow\text{exact coefficient-level moment integration}
\\
&\quad\longrightarrow\text{add the integrated interval remainder}.
\end{aligned}
}
\]
In particular, one should avoid replacing the polynomial by a range interval and multiplying that interval by the volume. That procedure discards cancellation and the exact moment structure of the polynomial.

The final enclosure has the form
\[
    \boxed{
    \int_{\mathcal{X}}f(x)\,dx
    \in
    I_{\mathrm{poly}}+R_{\mathrm{int}}[-1,1],
    }
\]
with
\[
    R_{\mathrm{int}}=E\abs{\mathcal{X}}
\]
for a constant pointwise remainder radius, or more generally
\[
    R_{\mathrm{int}}=\int_{\mathcal{X}}\rho(x)\,dx
\]
for a spatially varying certified remainder radius.

\end{document}
